\chapter{Introduction} \label{chapter:introduction}

Material interfaces are critical to the function and design of advanced solid-state systems. In semiconductors, they govern charge transport, carrier confinement, defect formation, and phase stability. In energy devices, they regulate ion exchange, chemical compatibility, and thermal response. In emerging heterostructures, interfaces can even induce novel behaviours not present in the constituent materials; such as encapsulated metastable phases or Schottky barriers. As device dimensions shrink and surface-to-volume ratios rise, the influence of interfaces grows, placing renewed emphasis on their accurate prediction and control.

Yet despite their importance, predicting which interfaces are energetically favourable remains a significant challenge. Interface stability is not solely determined by bulk lattice matching, but also by local bonding, chemical inhomogeneity, and atomic-scale reconstruction. The structural space is vast: each pair of surfaces can be stacked with arbitrary lateral shifts, terminations, and alignments, resulting in a combinatorial explosion of configurations. Strain relaxation, interlayer interactions, and charge transfer further complicate the energy landscape. Accurately assessing favourability across this space demands a balance between resolution and scale that conventional methods struggle to achieve.

First-principles techniques such as Density Functional Theory (DFT) remain the standard for computing accurate interfacial energies and relaxations. However, their high computational cost, scaling cubically with system size, renders exhaustive sampling intractable for realistic sizes structures, such as a 22 \times Si to 21 \times Ge latticed matched supercell, which often comprise hundreds of atoms. This bottleneck precludes broad screening efforts, especially when exploring multiple chemistries, orientations, and stacking arrangements. As such, DFT might be better suited to final validation and a confirmation of datasets evaluated in a less computationally expensive manner.

To address this, the present project explores the use of machine-learned interatomic potentials (MLPs) for scalable interface screening. MLPs are data-driven models trained to reproduce DFT-level energies and forces, but at orders of magnitude lower cost. In particular, this work focuses on the MACE framework; an equivariant message-passing architecture that preserves rotational symmetries and captures local atomic interactions with high fidelity. By using MACE to rapidly evaluate large numbers of candidate interfaces, it becomes possible to rank structures by predicted stability and prioritise low energetic configurations for subsequent DFT validation.

This approach supports a hierarchical modelling workflow in which DFT and MLPs are used in tandem: MLPs for breadth, and DFT for depth. It enables rigorous testing of whether MLPs can replicate DFT-derived favourability rankings across a diverse set of semiconductor systems, including phase boundaries (e.g. Si|Si) and heterostructures (e.g. Si|Ge, Ge|C). For the context of this study I ignore the semicontuctor requirement in its the definition. Moreover, it lays the foundation for further improvements through techniques such as delta-learning, structure-based prediction models, and surface-to-interface generalisation.

The remainder of this report is structured as follows:

\begin{itemize}
    \item \textbf{Chapter~\ref{chapter:background}} introduces the theoretical foundations and relevant literature, covering interface energetics, DFT formalism, and the development of MLPs.
    \item \textbf{Chapter~\ref{chapter:computational_investigation}} presents a benchmarking study comparing MLP and DFT predictions of interface favourability, including statistical correlation analysis and system-specific case studies.
    \item \textbf{Chapter~\ref{chapter:next_steps}} outlines planned extensions to the current methodology, including the development of delta-corrected potentials and predictive models linking surface properties to interface behaviour.
    \item \textbf{Chapter~\ref{chapter:conclusions}} concludes with a discussion of findings, limitations, and broader implications for autonomous materials modelling and the design of functional heterostructures.
\end{itemize}

Through this integrated modelling framework, the project seeks to reduce the cost and uncertainty of interface prediction while preserving the fidelity needed for scientific insight and technological relevance.